%%%%%%%%%%%%%%%%%%%%%%%%%%%%%%%%%%%%%%%%%%%%%%%%%%%%%%%%%%%%%%%%%%%%%%
% Lớp 10 | Phương pháp tọa độ trong mặt phẳng (𝑂𝑥𝑦) | Ba đường conic trong mặt phẳng tọa độ
% Dạng 0H9?5-4: Xác định các yếu tố của hypebol
% Số câu: 11  (sinh tự động bởi tach_cau.py; từ nay thầy sửa trực tiếp file này)
%%%%%%%%%%%%%%%%%%%%%%%%%%%%%%%%%%%%%%%%%%%%%%%%%%%%%%%%%%%%%%%%%%%%%%

%%%%%%%%%%%%%%%% TRẮC NGHIỆM 4 PHƯƠNG ÁN %%%%%%%%%%%%%%%%

% nguồn: MaHoaNH_TracNghiem_Lop10
\begin{ex}%[0H9N5-4]%[ma:1210b06e]
	Trong mặt phẳng $Oxy$, tọa độ các tiêu điểm của hypebol $(H)\colon\dfrac{x^2}{9}-\dfrac {y^2}{4}=1$ là 
	\choice 
		{\True $F_1(-\sqrt  {13},0)$; $F_2(\sqrt {13},0)$}
		{$F_1(0,-\sqrt {13})$; $F_2(0,\sqrt {13})$} 
		{$F_1(-\sqrt  {5},0)$; $F_2(\sqrt {5},0)$}
		{$F_1(0,-\sqrt {5})$;  $F_2(0,\sqrt {5})$}
	\loigiai{Ta  có $\heva{&a^2=9\\&b^2=4}$ khi đó $c^2=a^2+b^2=13\Rightarrow c=\sqrt{13}$. Vậy $F_1(-\sqrt {13},0)$; $F_2(\sqrt {13},0)$.}
\end{ex}

% nguồn: MaHoaNH_TracNghiem_Lop10
\begin{ex}%[0H9N5-4]%[ma:87cc09f5]
	Trong mặt phẳng tọa độ $Oxy$, cho hypebol $(H)\colon \dfrac{x^2}{16}-\dfrac{y^2}{9}=1$.  Độ dài trục ảo của hypebol $(H)$ bằng
	 \choice
		 {$8$}
		 {$10$}
		 {\True $6$}
		 {$12$}
	\loigiai{ 
		Từ phương trình chính tắc của  $(H)$, ta có $b^2=9\Rightarrow b=3$.\\
	Vậy độ dài trục ảo của $(H)$  là $2b=6$.}
\end{ex}

% nguồn: MaHoaNH_TracNghiem_Lop10
\begin{ex}%[0H9N5-4]%[ma:687cde81]
	 Xác định tiêu điểm của hypebol có phương trình $ (H)\colon \dfrac{x^2}{5}-\dfrac{y^2}{9}=1 $
	\choice 
		{$ F_2(2;0) $,  $ F_1(-2;0) $}
		 {$ F_2(0;\sqrt{14}) $}
		{$ F_1(0;-\sqrt{14}) $} 
		{\True $ F_2(\sqrt{14};0) $, $ F_1(-\sqrt{14};0) $} 
	 \loigiai{
		 Ta có $ a^2=5 $ và $ b^2=9 $ từ đó suy ra $ c=\sqrt{a^2+b^2}=\sqrt{5+9}=\sqrt{14} $.\\
		Vậy tiêu điểm của hybebol là $ F_2(\sqrt{14};0) $, $ F_1(-\sqrt{14};0)  $.
	} 
\end{ex}


%%%%%%%%%%%%%%%% ĐÚNG/SAI %%%%%%%%%%%%%%%%

% nguồn: MaHoaNH_DungSai_Lop10
\begin{ex}%[0H9N5-4]%[ma:ac3fde80]
	Cho hypebol có phương trình chính tắc $\dfrac{x^2}{9}-\dfrac{y^2}{16}=1$. Trong các mệnh  đề sau, mệnh đề nào đúng, mệnh đề nào sai?
	 \choiceTF
		{\True  Hypebol có tiêu điểm là $F_1(-5;0)$}
		 {Hypebol có tiêu điểm là $F_2(0;5)$}
		{\True Tiêu cự của hypebol là $10$} 
		{\True Hiệu các khoảng cách từ mỗi điểm trên hypebol tới hai tiêu điểm có giá trị tuyệt đối  bằng $6$}
	\loigiai{Ta có $a^2=9$, $b^2=16$.  Do đó, $c=\sqrt{a^2+b^2}=5$.\\
		Vậy $F_1(-5;0)$, $F_2(5;0)$ và tiêu  cự $F_1F_2=2c=10$.\\
		Từ đó, suy ra hiệu các khoảng cách từ mỗi điểm trên hypebol tới hai tiêu điểm  có giá trị tuyệt đối bằng $2a=6$.
	 }
\end{ex}

% nguồn: MaHoaNH_DungSai_Lop10
\begin{ex}%[0H9N5-4]%[ma:116caa55]
	Cho hypebol có phương trình chính tắc $\dfrac{x^2}{144}-\dfrac{y^2}{25}=1$. Trong các mệnh đề sau, mệnh  đề nào đúng, mệnh đề nào sai?
	\choiceTF 
		{\True Hypebol có tiêu điểm là $F_1(-13;0)$} 
		{Hypebol có  tiêu điểm là $F_2(0;13)$}
		{\True Tiêu cự  của hypebol là $26$}
		{\True Hiệu  các khoảng cách từ mỗi điểm trên hypebol tới hai tiêu điểm có giá trị tuyệt đối bằng $24$}
	\loigiai{Ta có $a^2=144$, $b^2=25$. Do đó, $c=\sqrt{a^2+b^2}=13$.\\ 
		Vậy $F_1(-13;0)$, $F_2(13;0)$ và  tiêu cự $F_1F_2=2c=26$.\\
		Từ đó,  suy ra hiệu các khoảng cách từ mỗi điểm trên hypebol tới hai tiêu điểm có giá trị tuyệt đối bằng $2a=24$.
	} 
\end{ex}

% nguồn: MaHoaNH_DungSai_Lop10
\begin{ex}%[0H9H5-4]%[ma:88019186]
	Cho hypebol có  dạng $\dfrac{x^2}{a^2}-\dfrac{y^2}{b^2}=1$ ($a,b>0$), đi qua điểm $A\left(\sqrt{3};0\right)$ và có một tiêu điểm $F_1(-2;0)$. Khi đó, các khẳng định sau là đúng hay sai?
	\choiceTF 
		 {Tiêu cự bằng $2$}
		{\True $  a=\sqrt{3}$}
		{$b^2=2  $}
		 {Điểm $B(0;1)$ thuộc hypebol}
	 \loigiai{
		\begin{itemchoice}
			\itemch \textbf{Sai}.\\ Hypebol có một tiêu  điểm $F_1(-2;0)$ nên $c=2$ và tiêu cự là $2c=4$.
			 \itemch \textbf{Đúng}.\\ Có $A\left(\sqrt{3};0\right)$ thuộc hypebol nên $\dfrac{\left(\sqrt{3}\right)^2}{a^2}-\dfrac{0^2}{b^2}=1\Rightarrow a^2=3\Rightarrow a=\sqrt{3}$.
			\itemch \textbf{Sai}.\\ Vì  $a=\sqrt{3}$, $c=2$ nên $b^2=c^2-a^2=1$.
			\itemch \textbf{Sai}.\\ Hypebol có phương trình  là $\dfrac{x^2}{3}-y^2=1$ không đi qua $B(0;1)$.
		\end{itemchoice}
	 }
\end{ex}

% nguồn: MaHoaNH_DungSai_Lop10
\begin{ex}%[0H9H5-4]%[ma:eb6a07dd]
	Trong mặt phẳng tọa độ, cho Hypebol $(H)$ đi qua điểm  $M\left(4; \sqrt{5}\right)$, có một tiêu điểm $F_1(-3;0)$ và đường thẳng $d\colon y=\dfrac{1}{2} x$
	 \choiceTF
		{\True Tiêu cự  của $(H)$ bằng $6$}
		 {\True $\left|MF_2\right|=\sqrt{6}$}
		{Gọi $A$ là một điểm tùy ý thuộc $(H)$. Khi đó $A$ cách trục  $Oy$ một khoảng ngắn nhất bằng $4$}
		{\True Đường thẳng $d$ cắt $(H)$ tại hai điểm  $E$, $F$ và $EF=2\sqrt{15}$}
	 \loigiai{
		\begin{itemchoice}
			\itemch Vì $(H)$ có tiêu điểm $F_1(-3; 0)$ nên $c=3$ suy  ra tiêu cự $2c=6$.
			\itemch  $(H)$ có tiêu điểm $F_2=(3; 0)$ nên $\left|MF_2\right|=\sqrt{(3-4)^2+\left(0-\sqrt{5}\right)^2}=\sqrt{6}$.
			\itemch Ta  có $c=3$ nên $\mathrm{d}(A;Oy)<3$.
			\itemch Gọi phương trình của  $(H): \dfrac{x^2}{a^2}-\dfrac{y^2}{b^2}=1$.\\
			Vì $M\left(4;\sqrt{5}\right) \in(H)$ nên ta  được: $\dfrac{16}{a^2}-\dfrac{5}{b^2}=1\quad\quad(1)$.\\
			Ta có $a^2+b^2=c^2$ mà $c=3$ nên  suy ra $b^2=9-a^2\quad\quad(2).$\\
			Thay $(2)$ vào $(1)$ ta được $\dfrac{16}{a^2}-\dfrac{5}{9-a^2}=1\Leftrightarrow a^4-30a^2+144=0\Leftrightarrow\left[\begin{array}{l}a^2=24\\ a^2=6.
		\end{array}\right.$\\
		 Mà $a^2< 9$ nên suy ra $a^2=6\Rightarrow b^2=3$.\\
		Do đó  $(H): \dfrac{x^2}{6}-\dfrac{y^2}{3}=1$.\\
		Tọa độ giao điểm của $d$ và $(H)$ là nghiệm  của hệ phương trình \allowdisplaybreaks
		\begin{eqnarray*}
			 &&\heva{&\dfrac{x^2}{6}-\dfrac{y^2}{3}=1\\ &y=\dfrac{1}{2} x}\\ &\Leftrightarrow&\heva{&\dfrac{4y^2}{6}-\dfrac{y^2}{3}=1\\ &x=2y}\\
			&\Leftrightarrow&\heva{& y^2=3 \\  & x=2y}\\
			&\Leftrightarrow&\heva{& y=\pm\sqrt{3} \\  & x=2y}
		\end{eqnarray*}
		Suy ra  $E\left(-2\sqrt{3};-\sqrt{3}\right), F\left(2\sqrt{3}; \sqrt{3}\right)$, $EF=\sqrt{\left(4\sqrt{3}\right)^2+\left(2\sqrt{3}\right)^2}=2\sqrt{15}$.
\end{itemchoice}}
\end{ex}

% nguồn: MaHoaNH_DungSai_Lop10
\begin{ex}%[0H9H5-4]%[ma:1a17025f]
	Trong  mặt phẳng $Oxy$, cho Hypebol $ (H) \colon \dfrac{x^2}{4}-\dfrac{y^2}{16}=1$.
	 \choiceTF
		{Điểm $A(0; 4)$ nằm  trên $(H)$}
		{\True Khoảng cách ngắn nhất giữa hai điểm nằm trên hai nhánh  của $(H)$ là $4$}
		{\True Hypebol $(H)$ có một tiêu điểm là $F_2\left( 2\sqrt{5};0\right)$} 
		 {\True Đường thẳng $d \colon x-y+1=0$ cắt $(H)$ tại hai điểm phân biệt}
	\loigiai{ 
	Gọi $2a$, $2b$ lần lượt là độ  dài trục thực, độ dài trục ảo của Hypebol đã cho.\\
	Ta có  $\heva{&a^2=4\\ &b^2=16} \Leftrightarrow \heva{&a=2\\ &b=4}$ và $c^2=a^2+b^2=20 \Rightarrow c=2\sqrt{5}$.
	\begin{itemize}
		\item Điểm $A(0;4)$ không  nằm trên Hypebol vì $\dfrac{0^2}{4}- \dfrac{4^2}{16}=-1 \neq 1$.
		\item Khoảng cách ngắn nhất giữa hai điểm nằm trên hai nhánh của $(H)$ là $A_1A_2=2a=2\cdot 2=4$. 
		\item Hypebol $(H)$ có một tiêu điểm là  $F_2\left( 2\sqrt{5};0\right)$ (vì $c=2\sqrt{5}$).
		\item Toạ  độ giao điểm của Hypebol $(H)$ và đường thẳng $d \colon x-y+1=0$ là nghiệm của hệ phương trình
		\[ 
		\heva{&\dfrac{x^2}{4}-\dfrac{y^2}{16}=1 \\ &x-y+1=0} \Leftrightarrow \heva{&4x^2-y^2=16  &(1)\\ &x-y+1=0.&(2)}
		\] 
		Từ $(2)$ suy ra $y=x+1$. Thay vào $(1)$ ta được 
		\[ 
		4x^2-  (x+1)^2=16 \Leftrightarrow 3x^2-2x-17=0 \Leftrightarrow x=\dfrac{1\pm 2\sqrt{13}}{3}.
		\] 
		Do đó đường thẳng $d$ cắt  hypebol $(H)$ tại hai điểm phân biệt.
	\end{itemize}
	\begin{itemchoice}
		 \itemch Sai.
		 \itemch Đúng.
		 \itemch Đúng.
		 \itemch Đúng.
	\end{itemchoice}
	} 
\end{ex}

% nguồn: MaHoaNH_DungSai_Lop10
\begin{ex}%[0H9V5-4]%[ma:8c5b53c1]
	Trong mặt phẳng tọa độ, cho hypebol $(H)$ có phương trình $9x^2-7y^2=63$  và đường thẳng $d\colon y=k x$.
	 \choiceTF
		 {\True $b^2=9$}
		{\True Giá trị tuyệt đối của hiệu các khoảng cách  từ mỗi điểm thuộc hypebol đến hai tiêu điểm bằng $2\sqrt{7}$}
		{\True Hai tiêu điểm là  $F_1(-4; 0), F_2(4; 0)$}
		{Có  $5$ giá trị nguyên của $k$ để đường thẳng $d$ và $(H)$ có điểm chung}
	 \loigiai{
		\begin{itemchoice}
			\itemch $9x^2-7y^2=63\Leftrightarrow  \dfrac{x^2}{7}-\dfrac{y^2}{9}=1$ nên suy ra $b^2=9$.
			\itemch $a^2=7\Leftrightarrow  a=\sqrt{7}$ suy ra $\left|MF_1-MF_2\right|=2a=2\sqrt{7}$.
			 \itemch $c^2=a^2+b^2=7+9=16\Rightarrow c=4$. Do đó hypebol có hai tiêu điểm $F_1(-4; 0), F_2(4; 0)$.
			\itemch Tọa độ giao điểm của $d$  và $(H)$ là nghiệm của hệ phương trình \allowdisplaybreaks
			\begin{eqnarray*}
				 &&\heva{&9x^2-7y^2=1\\&y=kx}\\
				&\Rightarrow& 9x^2-7. k^2x^2=1\\ 
				&\Leftrightarrow&  x^2\left(9-7k^2\right)=1\\
				&\Leftrightarrow& x^2=\dfrac{1}{9-7k^2}\quad(1). 
			\end{eqnarray*}
			Suy ra $d$ và $(H)$ có giao điểm thì phương trình $(1)$  có nghiệm
			 $$\Leftrightarrow 9-7k^2> 0\Leftrightarrow k^2< \dfrac{9}{7} \Leftrightarrow \dfrac{-3}{\sqrt{7}} < k < \dfrac{3}{\sqrt{7}}.$$
			Vậy $k$ nhận các giá trị nguyên $-1; 0; 1$ nên có ba  giá trị của $k$.
	\end{itemchoice}}
\end{ex}

% nguồn: MaHoaNH_DungSai_Lop10
\begin{ex}%[0H9V5-4]%[ma:82b28a9a]
	Trong mặt phẳng toạ độ, cho hypebol $(H)$ có  phương trình: $\dfrac{x^2}{9}-\dfrac{y^2}{16}=1$.
	\choiceTF 
		 {$a=4$}
		{\True Điểm $M(-3;0)$ nằm trên hypebol} 
		{\True Tiêu cự  của $(H)$ bằng $10$}
		{Điểm $M(3a;2a) \in(H)$ thuộc góc phần tư thứ nhất. Diện tích  tam giác $MF_1 F_2$ bằng $\dfrac{10\sqrt{3}}{3}$}
	\loigiai{ 
		\begin{itemchoice}
			\itemch  $a^2=9$ nên $a=3$.
			 \itemch Thay $x=-3$, $y=0$ vào phương trình $\dfrac{x^2}{9}-\dfrac{y^2}{16}=1$ được \allowdisplaybreaks
			\begin{eqnarray*}
				\dfrac{(-3)^2}{9}-\dfrac0{16}&=&1\\ 
				\Leftrightarrow  1&=&1\;\left(\text{đúng}\right)
			\end{eqnarray*}
			Do đó  $M\in(H)$.
			\itemch Ta có $a=3, b=4\Rightarrow c=\sqrt{a^2+b^2}=5\Rightarrow  2c=10$.\\
			Do đó tiêu cự  $F_1 F_2=2c=10$.
			\itemch Vì $M(3a;  2a) \in(H)$ nên ta được \allowdisplaybreaks
			\begin{eqnarray*}
				&&\dfrac{(3a)^2}{9}-\dfrac{(2a)^2}{16}=1\\ 
				&\Leftrightarrow&  a^2-\dfrac{a^2}{4}=1\Leftrightarrow a^2=\dfrac{4}{3}\\
				&\Leftrightarrow& a=\pm  \dfrac{2\sqrt{3}}{3}.
			\end{eqnarray*}
			Mà $M$ nằm trong  góc phần tư thứ nhất nên suy ra $M\left(2\sqrt{3}; \dfrac{4\sqrt{3}}{3}\right)$.\\
			 Diện tích tam giác $MF_1 F_2$ bằng: $S=\dfrac{1}{2} \cdot F_1 F_2 \cdot \dfrac{4\sqrt{3}}{3}=\dfrac{1}{2} \cdot 10\cdot \dfrac{4\sqrt{3}}{3}=\dfrac{20\sqrt{3}}{3}$.
	\end{itemchoice}}
\end{ex}


%%%%%%%%%%%%%%%% TỰ LUẬN %%%%%%%%%%%%%%%%

% nguồn: MaHoaNH_TuLuan_Lop10
\begin{bt}%[0H9H5-4]%[ma:7b5989cf]
	Trong mặt phẳng tọa độ $Oxy$ cho Hypebol $(H)$ có phương trình chính tắc là  $\dfrac{x^2}{6}-\dfrac{y^2}{8}=1$. Trên $(H)$ lấy điểm $M$ bất kì. Tính $\left|MF_1-MF_2\right|$.
	 \loigiai{Ta có $\heva{&a^2=6\\ &b^2=8} \Leftrightarrow\heva{&a=\sqrt{6} \\ &b=2\sqrt{2}}$ nên độ dài trục thực $2a=2\sqrt{6}$, độ dài trục ảo $2b=4\sqrt{2}$.\\
	Vì $M\in(H)$ nên ta có $\left|MF_1-MF_2\right|=2a=2\sqrt{6}$.} 
\end{bt}
